\section{官方幻灯片证据链：小样本评测为何仍需要统计学}

Sida Wang 的论证不是简单地说“小 benchmark 不可信”，而是把两个都合理的直觉放在一起检验：一方面，HumanEval、SWE-bench 等生成式或 Agent 评测远小于 MNIST/ImageNet；另一方面，一个真正困难且有长答案的问题似乎比一道选择题携带更多信息。整讲沿着逐题概率矩阵、配对比较和可预测噪声逐步回答：难题可以揭示新能力，但要稳定比较两个模型，仍然必须估计采样误差、推理随机性和题目分布。

\subsection{讲授路线与研究背景：先定义要测量的信号}

slides 1--3 给出题目、路线和讲者背景。这里的“millions of questions”不是单个 benchmark 有百万题，而是把许多模型、数据集、题目和重复采样组合成大规模逐题观测。讲者长期参与代码生成 benchmark 与 Eval-Arena，因此本讲关注的不是抽象统计练习，而是榜单上的几个百分点究竟能否支持研究决策。

\lecturefigure{slides-images/slide-001.png}{官方 slide 1：Predictable Noise and Patterns from Millions of Questions}
\lecturefigure{slides-images/slide-002.png}{官方 slide 2：从小型 benchmark 到噪声区间与改进建议的讲授路线}
\lecturefigure{slides-images/slide-003.png}{官方 slide 3：讲者在代码模型、Agent RL 与 benchmark 上的研究背景}

\begin{knowledgebox}{读图：先区分 capability signal 与 ranking signal}
某模型首次解决一个此前无人解决的问题，是强烈的 capability signal；但“模型 A 比模型 B 高 1.5 分”属于 ranking signal，需要重复样本和误差区间。前者可由一个高价值案例触发，后者必须回答差异是否能在新题目、新 seed 或新运行中复现。本讲主要解决第二类测量问题，同时保留第一类发现价值。
\end{knowledgebox}

\begin{knowledgebox}{术语消化：本讲的四类统计对象}
\begin{tabular}{p{3.0cm}p{4.2cm}p{5.0cm}}
\toprule
\textbf{术语} & \textbf{回答的问题} & \textbf{在本讲中的作用} \\
\midrule
Effect size & 两模型差多少 & 区分“有差异”与“差异足够有价值” \\
Standard error & 测得的差异会波动多少 & 决定当前测试集能否稳定分辨模型 \\
Confidence interval & 抽样程序给出的合理范围 & 防止把单点分数当成确定真值 \\
Signal-to-noise & 增益相对测量误差有多大 & 判断 benchmark 是否能测出目标级别的进步 \\
\bottomrule
\end{tabular}
\end{knowledgebox}

\subsection{样本量危机：难题是否能抵消测试集变小}

前一节区分了能力发现与精确排序，本节把差异落实到样本量与单题信息量。阅读接下来的规模对照时，需要同时追踪测试题数量、答案形态和研究声明：它们共同决定证据强度，任何单一维度都不足以宣布 benchmark 可靠或失效。

slides 4--8 先用历史对照建立压力。ImageNet 上 10\% 的提升足以改变领域判断，因为 10 万测试样本使统计显著性几乎不是争议；现代代码和 Agent benchmark 常只有数百题，每题却生成长代码、调用工具或执行多步轨迹。由此出现 A/B 两种假设：小数据集天然不可靠，或单个极难问题已经足够有信息量。

\lecturefigure{slides-images/slide-004.png}{官方 slide 4：ImageNet 式大提升与当前 LLM 小幅增益的可信度差异}
\lecturefigure{slides-images/slide-005.png}{官方 slide 5：MNIST、ImageNet 与生成式/Agent benchmark 的测试集规模}
\lecturefigure{slides-images/slide-006.png}{官方 slide 6：长代码答案是否让单题更有信息量}
\lecturefigure{slides-images/slide-007.png}{官方 slide 7：解决一个开放问题为何具有样本量为一的重大意义}
\lecturefigure{slides-images/slide-008.png}{官方 slide 8：小样本不可靠与少量极难题足够有效的两种竞争假设}

\begin{importantbox}{统计问题的准确问法}
测试集小并不自动意味着无用，答案长也不自动等于信息量大。需要问的是：观测结果对潜在能力变量有多敏感、重复运行的一致性多高、模型间差异相对标准误有多大。只有把“题目难度”和“响应可靠性”分开，才能判断一题究竟贡献了多少可比较信息。
\end{importantbox}

\textbf{课堂提示：}讲者拿 IMO、IOI 和开放问题作类比，不是要取消统计检验，而是提醒“发现能力”和“精确排序”具有不同证据门槛。研究报告应明确自己声称的是哪一种。

\subsection{逐题概率热图：强模型也会错简单题}

slides 9--12 把每列设为模型、每行设为按难度排序的问题，颜色表示同一模型在同一题上多次采样的 pass@1。热图并没有形成从左上到右下的清晰能力边界：弱模型偶尔解出难题，强模型也会在简单题失败。若单题代表稳定能力，这些反转不应如此常见；现实更像每个模型--题目组合都有一个连续成功概率。

\lecturefigure{slides-images/slide-009.png}{官方 slide 9：逐题、逐模型 pass probability 热图的阅读方法}
\lecturefigure{slides-images/slide-010.png}{官方 slide 10：不同代码 benchmark 的逐题概率矩阵之一}
\lecturefigure{slides-images/slide-011.png}{官方 slide 11：不同代码 benchmark 的逐题概率矩阵之二}
\lecturefigure{slides-images/slide-012.png}{官方 slide 12：HumanEval 中强弱模型跨难度反转的具体位置}

\begin{knowledgebox}{读图：颜色不是确定的 0/1 标签}
先横向比较同一题上的模型，再纵向比较同一模型跨题表现。大量中间颜色说明“会不会做”不是固定标签，而是一次采样的 Bernoulli 结果，其参数是潜在成功概率 $p_{m,i}$。难度排序只能解释平均趋势，不能保证每次运行都单调；因此单次 pass/fail 会丢掉概率信息。
\end{knowledgebox}

\begin{warningbox}{热图不证明所有题都一样}
不一致性可能来自模型随机性、题目歧义、测试缺陷、记忆污染或真实的能力组合差异。热图证明的是简单的一维“模型能力减题目难度”不足以解释全部数据，而不是证明题目内容无关。后续统计模型必须保留这种证据边界。
\end{warningbox}

\subsection{长答案并不自动提高统计功效}

slides 13--16 继续反驳“难题天然高信息量”。如果模型能写出看似深奥的证明，却无法定义基本概念，长答案可能来自复制、偶然搜索或不一致推理。slide 14 的关键观察是红色单元格并非 sharp zero：弱模型不是永远不会，只是成功概率更低。对可验证任务，首次从 0 提升到可重复的 1\% 也许揭示新能力，但从 1\% 到 99\% 仍需大量样本才能估计。

\lecturefigure{slides-images/slide-013.png}{官方 slide 13：用 hard question 补偿低统计功效的条件与反例}
\lecturefigure{slides-images/slide-014.png}{官方 slide 14：弱模型在难题上的非零概率与第一批成功样本}
\lecturefigure{slides-images/slide-015.png}{官方 slide 15：课堂重新选择 A/B 假设并否定无条件的小而难路线}
\lecturefigure{slides-images/slide-016.png}{官方 slide 16：把 kettle 误认成扫地机器人后的信任崩塌}

\begin{knowledgebox}{课堂提示：Kettle 例子为何重要}
在约 00:15:15，讲者描述维修助手先把烧水壶认成机器人吸尘器。即便后续建议碰巧正确，用户也无法确认中间推理是否可靠。多步 Agent 任务的总成功率近似由各关键步骤成功概率相乘；步骤越多，局部不一致越容易压低整条轨迹的可信度。长答案因此可能携带更多错误机会，而非更多独立信息。
\end{knowledgebox}

\subsection{Super-Population：把 benchmark score 放回抽样模型}

slides 17--20 转入统计基础。设题目 $x_i$ 从潜在任务总体 $\mathcal D$ 抽取，模型 $A$ 在题目上的得分为 $A(x_i)$。样本均值估计总体期望，样本方差估计题目间差异，均值标准误随 $1/\sqrt N$ 缩小。slide 19 再问观测差异是否可能仅由机会造成，slide 20 用数值例子显示小数据集的噪声区间足以吞掉常见的几个百分点增益。

\lecturefigure{slides-images/slide-017.png}{官方 slide 17：长答案争论之后仍需进入统计分析}
\lecturefigure{slides-images/slide-018.png}{官方 slide 18：从 super-population 抽样估计均值、方差与标准误}
\lecturefigure{slides-images/slide-019.png}{官方 slide 19：模型差异是否可能由 chance 产生}
\lecturefigure{slides-images/slide-020.png}{官方 slide 20：小测试集上的噪声区间数值量级}

\[
\widehat\mu_A=\frac{1}{N}\sum_{i=1}^{N}A(x_i),\qquad
\widehat{\mathrm{Var}}(A)=\frac{1}{N-1}\sum_{i=1}^{N}\bigl(A(x_i)-\widehat\mu_A\bigr)^2,
\]
\[
\widehat{\mathrm{SE}}(\widehat\mu_A)
=\sqrt{\frac{\widehat{\mathrm{Var}}(A)}{N}}.
\]
其中 $N$ 是测试题数，$\widehat\mu_A$ 是样本平均分，$\widehat{\mathrm{Var}}(A)$ 衡量不同题目上的离散程度，$\widehat{\mathrm{SE}}$ 表示若重新抽取测试集，均值估计会波动多少。标准误不是模型“真实能力的方差”，而是当前测量程序的不确定度。

\begin{warningbox}{置信区间不是“真值落在里面的概率”}
频率学派置信区间描述重复抽样程序的覆盖率。对一次固定实验，更稳妥的表达是“按该程序构造的区间具有 95\% 长期覆盖率”，而不是把参数本身说成以 95\% 概率位于区间内。讲义后文用区间作发布门禁时，也应保留这一解释。
\end{warningbox}

\subsection{Paired Comparison：同题比较为何更有力}

slides 21--24 区分 paired 与 unpaired。配对比较让模型 A、B 回答同一批题，再对逐题差值 $D_i=A(x_i)-B(x_i)$ 求均值；如果两模型在题目难度上的波动高度相关，公共难度项会在差分中抵消。非配对比较让两模型面对不同题目，题目组成差异会额外进入方差。slides 22--23 还加入推理随机变量，说明题目抽样与同题重复生成是两层不同噪声。

\lecturefigure{slides-images/slide-021.png}{官方 slide 21：Paired 与 unpaired 比较的直觉和方差公式}
\lecturefigure{slides-images/slide-022.png}{官方 slide 22：利用模型可独立重复采样估计推理噪声}
\lecturefigure{slides-images/slide-023.png}{官方 slide 23：题目方差与推理方差到 standard error 的分解}
\lecturefigure{slides-images/slide-024.png}{官方 slide 24：配对模型差值的标准误计算}

\[
\widehat\Delta=\frac1N\sum_{i=1}^{N}\bigl(A(x_i)-B(x_i)\bigr),\qquad
\widehat{\mathrm{SE}}(\widehat\Delta)
=\sqrt{\frac{\widehat{\mathrm{Var}}(D)}{N}}.
\]
这里 $D_i$ 是同一题上的模型差值。由
\[
\mathrm{Var}(A-B)=\mathrm{Var}(A)+\mathrm{Var}(B)-2\mathrm{Cov}(A,B)
\]
可见，当 A、B 都在相同难题上下降时，协方差为正，配对差值方差会明显小于两个独立均值方差之和。这就是相同测试集不仅“公平”，还提高统计功效的原因。

\begin{knowledgebox}{老师强调：seed variance 不是全部噪声}
模型可以对同一题独立采样，这让我们能估计条件于题目的推理方差；但重新抽一批题目还会产生任务总体方差。约 00:24:55，讲者提醒后者可能更重要。只跑多个 seed 而固定一小批题，无法回答结果对新任务是否稳定。
\end{knowledgebox}

\subsection{从 Bernoulli 怪现象到 Bootstrap、Sign Test 与 Eval-Arena}

slides 25--28 用“每次以 0.8 概率正确”的模型说明单次 0/1 观测为何会出现反直觉排序，再用数组运算、bootstrap 和 sign test 得到近似一致的误差判断。Bootstrap 从经验分布有放回重采样，模拟重新从 super-population 取题；sign test 只看两模型结果不同的题，并检验胜负不对称是否可能由机会造成。Eval-Arena 将这些 pairwise tests 系统化到大量模型对。

\lecturefigure{slides-images/slide-025.png}{官方 slide 25：真实正确率 0.8 仍会产生反直觉单次结果}
\lecturefigure{slides-images/slide-026.png}{官方 slide 26：用简单数组公式计算 Bernoulli 方差}
\lecturefigure{slides-images/slide-027.png}{官方 slide 27：Bootstrap 与 sign test 给出相近判断}
\lecturefigure{slides-images/slide-028.png}{官方 slide 28：Eval-Arena 对模型对执行配对比较和统计检验}

\begin{importantbox}{Bootstrap 的最小实现语义}
每次从 $N$ 个题目索引中有放回抽取 $N$ 个索引，重新计算模型差值；重复许多次后，差值分布近似描述测试集重采样的不确定性。必须按题目为单位一起重采样 A、B 的结果，才能保持配对结构；分别重采样会人为破坏协方差并扩大方差。
\end{importantbox}

\textbf{实践经验：}约 00:28:21，讲者认可 bootstrap，但建议先理解简单数组公式。工具不是替代理解：若重采样单位、配对关系或随机层级设错，运行再多次也只会得到精确的错误答案。

\subsection{Predictable Noise：数据集大小、准确率与相关性}

slides 29--32 给出本讲标题中的“可预测”。多个 benchmark 上，单模型标准误 $\mathrm{SE}(A)$ 与模型差值标准误 $\mathrm{SE}(A-B)$ 大致同量级，对应模型表现相关性约为 0.5；噪声还随总体 accuracy 呈弧形变化，因为二元得分方差 $p(1-p)$ 在 $p=0.5$ 最大、接近 0 或 1 时变小。SWE-bench Verified 上的筛选增益必须放入该准确率区间的噪声带解读。

\lecturefigure{slides-images/slide-029.png}{官方 slide 29：不同 benchmark 的可预测噪声与约 0.5 相关性}
\lecturefigure{slides-images/slide-030.png}{官方 slide 30：噪声随 overall accuracy 变化的经验曲线之一}
\lecturefigure{slides-images/slide-031.png}{官方 slide 31：更多数据集上的 accuracy-dependent noise}
\lecturefigure{slides-images/slide-032.png}{官方 slide 32：SWEFixer 与 P2P filter 的 SWE-bench Verified 配对结果}

\begin{knowledgebox}{读图：为什么准确率中段最吵}
对二元结果，单次方差为 $p(1-p)$。当模型几乎总错或总对时，重复结果较稳定；当 $p\approx0.5$ 时最不稳定。不过模型比较看的是 $A-B$，还受两者逐题相关性影响。图中的经验曲线因此是 Bernoulli 基线、题目异质性和模型相关性的合成，不应只套一个理论抛物线。
\end{knowledgebox}

\subsection{Beta 概率形态与“重新加权题目”的失败}

slides 33--36 从逐题成功概率的经验分布继续建模。强模型的题目概率常向 0 和 1 两端集中，可近似呈双峰 Beta 形态；但该观察主要来自 binary predictions，不能直接外推到连续 judge score。讲者原本希望通过难度建模、过滤或重加权降低噪声，却发现 Item Response Model 等方法没有稳定优于朴素基线，主因是模型自身响应不一致，而非少数“坏题”独占噪声。

\lecturefigure{slides-images/slide-033.png}{官方 slide 33：逐题经验成功概率及 Beta 分布近似}
\lecturefigure{slides-images/slide-034.png}{官方 slide 34：不同模型和 benchmark 的 Beta 累积分布}
\lecturefigure{slides-images/slide-035.png}{官方 slide 35：二元预测下更多 Beta 经验曲线}
\lecturefigure{slides-images/slide-036.png}{官方 slide 36：过滤、重加权与题目难度建模未能显著提升 signal}

\begin{warningbox}{Beta 拟合的证据边界}
Beta 是 $[0,1]$ 概率的灵活分布，可描述两端集中，但拟合良好不等于生成机制就是 Beta，也不证明所有题目可由单一难度参数解释。对长文本 rubric、部分得分和工具轨迹，应先定义观测模型，再决定是否仍适用 Bernoulli/Beta 框架。
\end{warningbox}

\begin{importantbox}{课堂失败结论}
“清洗得更精致”不一定增加可比较信息。HumanEval+、MBPP+ 或主观更高质量的数据集未必有更高 signal-to-noise；当主要噪声来自模型随机性和能力不一致时，仅过滤题目会缩小 $N$，反而提高标准误。优先策略是增加 benchmark 数量、扩大样本，或让每个昂贵 Agent 样本产出超过 1 bit 的结构化信息。
\end{importantbox}

\subsection{测量建议、Signal-to-Noise 与 Agent Eval 治理}

slides 37--42 收束为可执行建议。multiple seeds 和训练曲线方差有用，但通常低估完整不确定性；更可靠的顺序是读取已有噪声表、按 accuracy 区间查曲线、导出逐题结果运行配对代码。signal-to-noise 定义为模型系列中有意义增益相对标准误的比值；许多代码生成 benchmark 低于 2，无法稳定测出模型规模翻倍的增益。最后的 SWE-bench issue 提醒：当一次评测涉及 10 万 token、无人审计轨迹时，数据和执行错误会长期隐藏。

\lecturefigure{slides-images/slide-037.png}{官方 slide 37：multiple seeds、训练曲线与逐题结果的噪声测量建议}
\lecturefigure{slides-images/slide-038.png}{官方 slide 38：Eval-Arena 的噪声参考、排名与可疑样本工具}
\lecturefigure{slides-images/slide-039.png}{官方 slide 39：各 benchmark 的 signal-to-noise 排名表}
\lecturefigure{slides-images/slide-040.png}{官方 slide 40：signal-to-noise 还依赖模型类别}
\lecturefigure{slides-images/slide-041.png}{官方 slide 41：复杂 SWE-bench 执行链路中的真实审计事故}
\lecturefigure{slides-images/slide-042.png}{官方 slide 42：不同模型对和规模区间的 paired comparison 结果}

\[
\mathrm{SNR}_{A,B}
=\frac{|\widehat\mu_A-\widehat\mu_B|}
{\widehat{\mathrm{SE}}(\widehat\mu_A-\widehat\mu_B)}.
\]
分子是要测量的模型增益，分母是同题配对差值的标准误。$\mathrm{SNR}<2$ 通常意味着差异尚不足以稳定越过常见 95\% 双侧阈值，但它不是跨所有多重比较情形的万能判据。排行榜同时比较大量模型时，还需控制 family-wise error 或 false discovery rate。

\begin{knowledgebox}{读图：Signal-to-noise 比“题目看起来难”更接近 benchmark 质量}
slide 39 中，MMLU、TQA 等大规模选择题的 SNR 高于许多复杂代码 benchmark；这不表示选择题更接近真实工作，而是表示它们更稳定地测出既定模型差异。生态有效性与统计可靠性是两条轴：Agent benchmark 可以更真实，却仍需要更多任务、更丰富轨迹标签和更严格执行审计。
\end{knowledgebox}

\begin{warningbox}{Agent Eval 的额外故障面}
Agent 分数不仅受题目和推理 seed 影响，还受环境镜像、依赖版本、工具权限、超时、重试和 evaluator 脚本影响。复杂度越高，越不能只保存最终 pass/fail。应保留逐步轨迹、工具返回、环境哈希和失败归因，使“模型没做对”与“评测基础设施坏了”能够分开。
\end{warningbox}

\subsection{本章小结}
官方 42 页幻灯片形成一条完整证据链：小 benchmark 的长答案并未自动解决统计功效；逐题概率揭示模型响应具有连续性和不一致性；paired comparison 与方差分解给出可计算的不确定度；经验噪声又能由数据集大小、accuracy 和模型相关性部分预测。最终建议不是停止评测，而是报告误差、共享逐题结果、扩大信息量并审计复杂执行链路。

